%% kapa: exam=1 round=22 year=2011 period=2 form=A subject=accounting subject_ko=회계학 seq=30 number=30
%% source: Q-Net 감정평가사 기출문제 (cst00309) artlSeq=5036790 file=extracted/1008/5036790/2011년 제22회 감정평가사 1차시험 기출문제/2교시.hwp
%% conversion: kordoc 4.24.51 -> tools/gen1.py (auto); review=proofread-sonnet; flags=none
\begin{question}{30}
(주)서울은 20×1년 1월 1일 자금조달을 목적으로 액면금액이 ￦500,000인 사채(액면이자 ￦40,000 매년 말 지급, 시장이자율 연 10\%, 3년 만기)를 ￦475,122에 발행하였으며, 동 사채를 동 일자에 (주)한국이 발행금액에 취득하고 만기보유금융자산으로 분류하여 회계처리하였다. (주)한국은 20×1년 12월 31일 보유목적 등의 변화로 인하여 동 만기보유금융자산을 매도가능금융자산으로 분류변경하였다. 한편 20×1년 말 시장이자율이 연 12\%로 상승하였으며, 20×1년 말과 20×2년 말 동 자산의 공정가치는 다음과 같다.

{\small\centering
\begin{tabular}{|l|r|r|}
\hline
 & \multicolumn{1}{c|}{\uline{20×1년 말}} & \multicolumn{1}{c|}{\uline{20×2년 말}} \\
\hline
공정가치 & ￦466,204 & ￦482,151 \\
\hline
\end{tabular}\par}

(주)한국이 동 자산과 관련하여 20×1년 포괄손익계산서에 인식할 평가손실과 20×2년 포괄손익계산서에 인식할 이자수익은? (단, 발행 및 취득과 관련하여 거래비용은 발생하지 않았다고 가정, 최초인식 및 분류변경 또한 적절하였다고 가정, 현가계수는 다음과 같음) \par
(2년 기준) 연 12\% \par
단일금액 ￦1의 현가계수 0.7972 \par
정상연금 ￦1의 현가계수 1.6901 \par
\noindent\hspace*{2em}\makebox[3.2cm][l]{\uline{20×1년 평가손실}}\uline{20×2년 이자수익}\par
\begin{choices}
  \item \makebox[3.2cm][l]{￦8,918}￦47,512
  \item \makebox[3.2cm][l]{￦13,253}￦47,512
  \item \makebox[3.2cm][l]{￦13,253}￦48,263
  \item \makebox[3.2cm][l]{￦16,430}￦48,263
  \item \makebox[3.2cm][l]{￦16,430}￦57,916
\end{choices}
\end{question}
\answer{4}
